\documentclass[11pt]{article}
\usepackage[letterpaper,margin=1in]{geometry}
\usepackage[T1]{fontenc}
\usepackage[utf8]{inputenc}
\usepackage{amsmath,amssymb,amsthm,lmodern,microtype,booktabs}
\usepackage[hyphens]{url}
\usepackage[hidelinks]{hyperref}

\newtheorem{theorem}{Theorem}
\newtheorem{proposition}[theorem]{Proposition}
\newtheorem{lemma}[theorem]{Lemma}
\newtheorem{corollary}[theorem]{Corollary}
\theoremstyle{remark}
\newtheorem{remark}[theorem]{Remark}

\title{Five-Cut Profiles and Near-Bipartite Five-Poles\\for the Berge--Fulkerson Conjecture}
\author{}
\date{}

\begin{document}
\maketitle

\begin{abstract}
We present an independent, DRAT-certified replay of the five-cut boundary-profile
classification of M\'{a}\v{c}ajov\'{a}, Mazzuoccolo, and Tabarelli~\cite{MMT} for a
minimum counterexample to the Berge--Fulkerson conjecture. Under the Kempe-switch and
acyclic-completion necessary conditions, we recover all $1{,}716$ ordered disjoint
profile pairs, forming $13$ orbits under simultaneous port relabeling and side exchange.
Twelve orbits match the published pentagon cases I--XII, while the thirteenth corrects
an exchange of the port-pair labels $25$ and $45$ in the published Figure~7 and Appendix
for case~XIII. We then determine exact boundary profiles for two explicit families of
five-poles: bipartite five-poles (via six ordinary matching tests) and five-poles with
five terminals on one side of a single internal edge defect. Under cyclic edge-connectivity
hypotheses, bipartite five-poles realize all $21$ bipartite types (and yield
$3$-edge-colorable gluings with pentagons and other bipartite five-poles), while the
one-edge-defect five-poles realize $36$, $45$, or $55$ types. Together with the
five-cut classification, both pole families are excluded from cyclic $5$-edge cuts of a
minimum counterexample.
\end{abstract}

\section{Preliminaries}

The Berge--Fulkerson conjecture~\cite{Fulkerson1971} asserts that every bridgeless cubic
graph admits six perfect matchings covering each edge exactly twice. Equivalently, one
assigns a $2$-element subset of six colors $\{0,1,2,3,4,5\}$ to each edge so that the
three pairs incident with every vertex partition the six colors.

A \emph{five-pole} $P$ is a cubic graph fragment with five dangling semiedges (ports),
labeled $1,\dots,5$, attached at distinct terminal vertices; internal non-terminal
vertices have degree~$3$ in~$P$ and terminals have degree~$2$ in~$P$. Since
$3|V(P)| = 2|E(P)| + 5$, the order $|V(P)|$ is odd. In any Berge--Fulkerson coloring of
$P$, each color $c$ has an odd-sized boundary support $S_c \subseteq \{1,2,3,4,5\}$, and
each port belongs to exactly two color supports, so $\sum_{c=0}^{5} |S_c| = 10$.
Forgetting color names by sorting the six supports gives the \emph{boundary type} of the
coloring. There are $56$ possible five-port boundary types:
\begin{itemize}
  \item the unique \emph{all-five type}, with one color using all five ports and the
        other five using one port each; and
  \item $55$ types with two \emph{heavy} colors using three ports each and four
        \emph{light} colors using one port each.
\end{itemize}
Following~\cite{MMT}, we denote a heavy type by $(ij)(kl)$, where $ij$ and $kl$ are the
unordered pairs of ports omitted by the two three-port colors; when $ij = kl$ the type
is a \emph{loop} $(ij)(ij)$. The four singleton multiplicities are uniquely forced by
the condition that every port appear in two colors. The \emph{profile} of a five-pole is
the set of boundary types realized by its Berge--Fulkerson colorings. Two five-poles
glued along matching ports form a cubic graph admitting a Berge--Fulkerson cover if and
only if their profiles intersect.

\section{Correcting the Port Labels in Case~XIII}

Proposition~4.4 of~\cite{MMT} classifies the possible disjoint profile pairs at a cyclic
$5$-edge cut of a minimum counterexample into thirteen orbits: cases I--XII pair an
actual pentagon ($16$ types) with a disjoint partner profile, while case~XIII pairs two
non-pentagon profiles.

\begin{proposition}\label{prop:case13-error}
The right-hand profile of case~XIII as printed in Figure~7 and the Appendix of~\cite{MMT}
cannot be the profile of any five-pole, under any ordering of the five ports.
\end{proposition}

\begin{proof}
The printed right-hand profile of case~XIII contains the all-five type. Consider any
Berge--Fulkerson coloring of a five-pole realizing the all-five type, and let $c_0$ be
the color using all five ports and $c_2$ be the singleton color using port~$2$. The
subgraph of internal edges and semiedges carrying exactly one of $c_0, c_2$ has degree
$0$ or $2$ at every internal vertex, so its open paths pair up the four exclusive
boundary ports $\{1,3,4,5\}$. The open path starting at port~$1$ ends at $3$, $4$, or
$5$. Swapping $c_0$ and $c_2$ along this path produces a valid coloring of boundary type
$(13)(45)$, $(14)(35)$, or $(15)(34)$, respectively. (Switching the complementary path
yields the same boundary type up to exchanging the names of $c_0$ and $c_2$.) Hence any
valid five-pole profile containing the all-five type must contain at least one of
$(13)(45)$, $(14)(35)$, and $(15)(34)$. The right-hand profile printed in~\cite{MMT}
omits all three.
\end{proof}

Our independent replay (Section~\ref{sec:replay}) recovers the unique valid thirteenth
orbit. Its left-hand profile ($30$ types) consists of the two looped cliques on the
complementary $5$-cycles in $K_5$:
\[
  A = \{13, 14, 24, 25, 35\}, \qquad B = \{12, 15, 23, 34, 45\}.
\]
Its right-hand profile ($26$ types) consists of all $25$ cross pairs in $A \times B$
together with the all-five type. In particular, $(13)(45)$ belongs to $A \times B$ and
hence appears in the corrected right-hand profile. Comparing with~\cite{MMT}, Figure~7
interchanges the labels $25$ and $45$ between $A$ and $B$. In the Appendix of~\cite{MMT},
the left-hand profile lists the two blocks in the order $A, B$ above and repeats $45$ in
both blocks; replacing $45$ by $25$ in the first block of the left-hand profile and
exchanging $25$ and $45$ in the right-hand profile restores the corrected case~XIII. In
the $56$-bit ordering of \texttt{five\_cut.py}, the corrected pair is
\texttt{b4abccc8f986b9} ($30$ types) and \texttt{4b543337067946} ($26$ types).

\section{Exact Replay and DRAT Exhaustion Certificate}\label{sec:replay}

Let $G$ be a minimum counterexample to the Berge--Fulkerson conjecture (minimum order
among simple bridgeless cubic graphs with no cover). By standard small-cut reductions
(see~\cite{MMT}), $G$ is cyclically $5$-edge-connected and has girth at least~$5$. At
any cyclic $5$-edge cut separating $G$ into two five-poles $P$ and $Q$, both poles
contain a cycle and therefore have at least $5$ vertices.

\begin{lemma}\label{lem:pole-connected}
In a cyclically $5$-edge-connected cubic graph, each side $P$ of a cyclic $5$-edge cut is
connected, bridgeless, and has five distinct terminal vertices.
\end{lemma}

\begin{proof}
If a vertex $v \in V(P)$ carried two cut edges, then $v$ would have internal degree~$1$
in $P$, so $P - v$ would still contain a cycle while moving $v$ to $Q$ would yield a
cyclic $4$-edge cut; similarly a vertex with three cut edges would yield a cyclic
$2$-edge cut. Thus all five terminals are distinct and every vertex of $P$ has internal
degree at least~$2$, so every component of $P$ contains a cycle. Two components would
each require at least $5$ cut edges (since each complement includes the cyclic pole~$Q$),
impossible with five cut edges total; hence $P$ is connected.

Suppose $P$ had an internal bridge $e$ separating $P$ into components $P_1, P_2$. If
$P_i$ were a tree on $n$ vertices, its boundary in $G$ would have size $n + 2$, whereas
distinct terminals give at most $n$ terminal edges plus the bridge $e$, a contradiction.
Hence both $P_1$ and $P_2$ contain cycles, so each requires at least $4$ terminals in
addition to $e$ to avoid a cyclic cut of size at most~$4$ against $Q$; this requires at
least $8$ terminals, a contradiction.
\end{proof}

We encode two families of necessary conditions satisfied by the disjoint profiles
$\Pi(P)$ and $\Pi(Q)$ of a cyclic $5$-edge cut in a minimum counterexample:

\begin{enumerate}
  \item \emph{Single and simultaneous Kempe switches.} For every boundary type
        $\alpha \in \Pi(P)$ and every pair of colors $\{c_1, c_2\}$, the exclusive ports
        in $S_{c_1} \triangle S_{c_2}$ have even size at most~$4$ and are paired by open
        Kempe chains. When there are $4$ exclusive ports, switching either of the two
        chains produces the same boundary type up to exchanging $c_1$ and $c_2$, giving
        $3$ possible switch outcomes, at least one of which must belong to $\Pi(P)$. For
        any two disjoint pairs of colors $\{c_1,c_2\}$ and $\{c_3,c_4\}$, swaps act on
        disjoint color bits and therefore commute, so $\Pi(P)$ must simultaneously
        contain both single-pair outcomes and the combined double-switch outcome for some
        choice of chain pairings.
  \item \emph{Acyclic completions.} Consider the two acyclic five-poles from Figure~4
        of~\cite{MMT}:
        \begin{itemize}
          \item a single cubic vertex incident with ports $1,2,3$ together with a bare
                edge joining ports $4$ and $5$; its profile is the looped clique on
                $\{12, 13, 23\}$ ($6$ types, $10$ orientations under $S_5$); and
          \item a three-vertex path $u\text{--}w\text{--}v$ with ports $1,2$ at $u$,
                port $3$ at $w$, and ports $4,5$ at $v$; its profile is the looped clique
                on $\{14, 15, 24, 25\}$ ($10$ types, $15$ orientations under $S_5$).
        \end{itemize}
        By Lemma~\ref{lem:pole-connected}, gluing either cap to $P$ creates only edges on
        cycles (since every branch of a cap reaches a terminal on the connected bridgeless
        pole~$P$) and reduces the vertex count (from at least $5$ vertices in $Q$ to $1$
        or $3$). If the bare edge of the first cap joins two terminals that are already
        adjacent in $P$, the resulting double edge is harmless: in any bridgeless cubic
        multigraph smaller than $G$, a triple-edge component trivially admits a cover,
        while suppressing a double edge and its two endpoints into a single edge between
        their two distinct external neighbors yields a still smaller bridgeless cubic
        multigraph whose cover lifts immediately. By minimality of $G$, every glued graph
        has a cover, so both $\Pi(P)$ and $\Pi(Q)$ intersect every orientation of both
        caps. Since $\Pi(P) \cap \Pi(Q) = \varnothing$, each profile must also omit at
        least one type from every orientation of both caps.
\end{enumerate}

\begin{theorem}\label{thm:replay}
Over all pairs of disjoint subsets of the $56$ five-port boundary types satisfying the
Kempe-switch rules ($430$ distinct implications) and the $25$ acyclic-completion
intersection and non-containment constraints, there are exactly $1{,}716$ ordered
disjoint pairs. Under simultaneous port relabeling and exchange of the two sides, they
form exactly $13$ orbits: cases I--XII of~\cite{MMT} (with partner profile sizes
$25, 28, 32, 32, 33, 33, 34, 36, 37, 38, 39, 40$ paired with the $16$-type pentagon) and
the corrected case~XIII ($30$ and $26$ types) of Section~2.
\end{theorem}

The SAT enumeration and independent proof check are implemented in
\path{replay_five_targets.py}. Adding one blocking clause for each of the $1{,}716$
archived pairs in \path{certificates/bf-five-profile-exhaustion.pairs} yields an
unsatisfiable CNF formula (\path{bf-five-profile-exhaustion.cnf}, $3{,}652$ variables
and $13{,}172$ clauses) with a \texttt{drat-trim}-verified DRAT refutation
(\path{bf-five-profile-exhaustion.drat}):
{\small\begin{verbatim}
python3 replay_five_targets.py --verify certificates/bf-five-profile-exhaustion
drat-trim certificates/bf-five-profile-exhaustion.cnf \
          certificates/bf-five-profile-exhaustion.drat
\end{verbatim}}

\section{Exact Profiles of Bipartite Five-Poles}

We now determine the exact profile of any bipartite five-pole and exclude it from cyclic
$5$-edge cuts of a minimum counterexample.

\begin{proposition}\label{prop:bipartite-six}
Let $P$ be a five-pole whose internal graph is bipartite with vertex bipartition $(A,B)$.
After interchanging $A$ and $B$ if necessary, $A$ contains four terminals, $B$ contains
one terminal $b$, and $|A| = |B| + 1$. A three-port mask is realizable by a single color
if and only if it consists of $b$ and two $A$-terminals $u, v$ such that $P - \{b,u,v\}$
admits a perfect matching. If $t \in \{0,\dots,6\}$ of these six triples are feasible,
then $\Pi(P)$ consists precisely of the $\binom{t+1}{2}$ heavy types formed by choosing
any multiset of two feasible triples.
\end{proposition}

\begin{proof}
Let $p_A, p_B$ be the numbers of terminals in $A, B$, so $p_A + p_B = 5$. Counting
internal edges from both sides gives $3|A| - p_A = 3|B| - p_B$, i.e.,
$3(|A| - |B|) = p_A - p_B$. Among odd integers of absolute value at most~$5$, only
$\pm 3$ is divisible by~$3$, so swapping $A$ and $B$ if necessary gives $p_A = 4$,
$p_B = 1$, and $|A| = |B| + 1$.

In any color $c$, let $s_A, s_B$ be the numbers of $A$- and $B$-terminals in $S_c$.
Internal edges of color $c$ form a matching between $A \setminus S_c$ and
$B \setminus S_c$, so $|A| - s_A = |B| - s_B$, hence $s_A - s_B = 1$. Since
$s_B \in \{0,1\}$, either $s_B = 0$ and $s_A = 1$ (a singleton in $A$), or $s_B = 1$ and
$s_A = 2$ (a triple $\{b, u, v\}$ with $u, v \in A$). Because port~$b$ carries two
colors, every Berge--Fulkerson coloring of $P$ has exactly two triple colors (both
containing $b$) and four singleton colors. Necessarily, each triple $\{b, u, v\}$ is
feasible only if $P - \{b, u, v\}$ has a perfect matching.

Conversely, choose any multiset of two feasible triples $T_1, T_2$ (possibly equal) and
fix perfect matchings $M_1, M_2$ of $P - T_1$ and $P - T_2$. Start with two copies of
every internal and dangling edge of $P$, and allocate one copy of $M_1 \cup T_1$ and one
copy of $M_2 \cup T_2$ to the first two colors. In the residual bipartite multigraph,
every vertex has degree~$4$ counting remaining dangling copies; none of the remaining
dangling copies lies on $B$, and exactly four lie on $A$. Attach those four dangling
copies to a single new vertex $x$ on the $B$-side to obtain a $4$-regular bipartite
multigraph (with $|A| = |B| + 1$ vertices on each side). By K\H{o}nig's/Hall's theorem,
its edge set decomposes into four perfect matchings, each using one edge incident to~$x$
and hence supplying the four required singleton colors.
\end{proof}

\begin{theorem}\label{thm:bipartite-21}
Suppose a bipartite five-pole $P$ is one side of a cyclic $5$-edge cut in a simple cubic
graph $G$ having no cyclic edge cut of size at most~$3$. Then:
\begin{enumerate}
  \item Every triple $\{b, u, v\}$ is feasible, so $|\Pi(P)| = \binom{7}{2} = 21$.
  \item $P$ shares exactly $6$ boundary types with a pentagon in any port ordering, and
        two such bipartite five-poles share $21$ types if their minority ports align and
        $6$ otherwise.
  \item Every gluing of $P$ with a pentagon or with another such bipartite five-pole is
        $3$-edge-colorable.
\end{enumerate}
\end{theorem}

\begin{proof}
Fix any two $A$-terminals $u, v$. We verify Hall's condition for $P - \{b, u, v\}$ from
$B \setminus \{b\}$ to $A \setminus \{u, v\}$. Let $Y \subseteq B \setminus \{b\}$ be
nonempty and set $N = N_P(Y) \subseteq A$. Every vertex of $Y$ is a non-terminal, so all
$3|Y|$ edges incident to $Y$ enter $N$. If $N$ contains $p$ terminals, then
$3|Y| \le 3|N| - p$, giving $3(|N| - |Y|) \ge p \ge 0$.

If Hall's condition fails on $Y$, then $|N \setminus \{u,v\}| < |Y|$, which requires
$|N| - |Y| \in \{0, 1\}$. If $|N| - |Y| = 0$, then $p = 0$, so $u, v \notin N$ and
$|N \setminus \{u,v\}| = |N| = |Y|$, a contradiction. Thus $|N| - |Y| = 1$ and both
$u, v \in N$. Since $G$ is simple, a single vertex in $Y$ has three distinct neighbors in
$N$, so $|Y| \ge 2$. In $G$, the set $N \cup Y$ has $2|Y| + 1$ vertices, $3|Y|$ induced
edges (strictly more than a tree, hence containing a cycle), and boundary size
$3|N| - 3|Y| = 3$. Its complement contains the other cyclic pole of the $5$-cut, yielding
a cyclic $3$-edge cut in $G$, a contradiction. Thus all $6$ triples are feasible, giving
$21$ types.

For part~(2), a pentagon's heavy triples omit adjacent cycle edges; removing port~$b$
from the $5$-cycle leaves a $4$-vertex path with $3$ edges, giving $\binom{4}{2} = 6$
common types. For two bipartite five-poles with distinct minority ports $b_1 \neq b_2$,
common triples omit pairs among the $3$ mutually majority ports, again giving
$\binom{4}{2} = 6$ types.

For part~(3), pick a common heavy triple $T$ and combine partial perfect matchings on
both poles with $T$ to form a perfect matching $M$ of the glued graph $H$. The
complement $H - M$ is $2$-regular and uses only the two cut edges $e_1, e_2$ outside~$T$,
which must therefore lie on a single cycle $C$ of $H - M$. On each bipartite side, $T$
contains the minority port, so $e_1$ and $e_2$ both meet majority vertices and the
internal path of $C$ joining them has even length. On a pentagon, deleting the three
vertices of $T$ leaves one matched edge in $M$, so the internal path of $C$ is a path of
length~$4$ (even) and no other component remains. Every other cycle of $H - M$ lies inside
a bipartite pole and is even. Hence $H - M$ is bipartite and decomposes into two perfect
matchings, giving a $3$-edge-coloring of $H$.
\end{proof}

\section{Five Terminals on One Side of a One-Edge Defect}

\begin{theorem}\label{thm:one-defect}
Suppose a five-pole $P$ is one side of a cyclic $5$-edge cut in a simple cyclically
$5$-edge-connected cubic graph $G$, and suppose $V(P)$ admits a partition $(A, B)$ with
all five terminals in $B$, no internal edge within $B$, and exactly one internal edge
$pq$ within $A$. Then $|\Pi(P)| \in \{36, 45, 55\}$, and $\Pi(P)$ shares at least $6$
boundary types with every orientation of a pentagon.
\end{theorem}

\begin{proof}
Counting crossing edges between $A$ and $B$ gives $3|A| - 2 = 3|B| - 5$, so
$|B| = |A| + 1$. For a single color $c$ using $s$ dangling edges (all on $B$) and
$x \in \{0,1\}$ copies of $pq$, matching crossing edges gives $|A| - 2x = |B| - s$, so
$s = 1 + 2x$. Thus the two colors on $pq$ use three ports each, and the other four colors
are singletons. A heavy triple (equivalently, its omitted pair of terminals $\{u, v\}$)
is feasible for a single color if and only if the bipartite graph $J = P - \{p, q\}$
admits a perfect matching between $A_J = A \setminus \{p, q\}$ and
$B_0 \cup \{u, v\}$, where $B_0$ is the set of non-terminal vertices of~$B$. Conversely,
given any multiset of two feasible omitted pairs, allocating their matchings and both
copies of $pq$ leaves a bipartite residual multigraph with four dangling edge copies on
$B$. Attaching them to one new vertex on $A$ yields a $4$-regular bipartite multigraph
whose decomposition into four perfect matchings supplies the four singletons. Thus if
$t$ omitted pairs are feasible, $|\Pi(P)| = \binom{t+1}{2}$.

For $W \subseteq B$, let $k$ be the number of terminals in $W$ and define the deficiency
$\delta(W) = |W| - |N_J(W)|$. Consider $S = W \cup N_J(W) \cup \{p, q\} \subseteq V(P)$.
The induced edges of $G[S]$ are precisely the $3|W| - k$ internal edges incident to $W$
together with $pq$, so the boundary size of $S$ in $G$ is
\[
  |\partial_G(S)| = 3(|W| + |N_J(W)| + 2) - 2(3|W| - k + 1) = 4 + 2k - 3\delta(W).
\]
If $W \neq \varnothing$, then $|S| \ge 3$. In a cubic graph, any forest on $n \ge 3$
vertices has boundary at least $n + 2 \ge 5$, so $|\partial_G(S)| \le 4$ would force a
cycle inside $S$. Since $V(G) \setminus S$ contains the complementary cyclic pole of the
$5$-cut, cyclic $5$-edge-connectivity forces $4 + 2k - 3\delta(W) \ge 5$ for every
nonempty $W \subseteq B$. Because $\delta(\varnothing) = 0$ and $\delta(W)$ is an
integer, it follows that $\delta(W) \le 0$ whenever $k \le 1$, and $\delta(W) \le 1$
whenever $k \le 3$.

If an omitted pair $\{u, v\}$ is infeasible, Hall's theorem gives a witness
$W_{uv} \subseteq B_0 \cup \{u, v\}$ with $\delta(W_{uv}) \ge 1$. Since $\delta \le 0$
for $k \le 1$, $W_{uv}$ must contain both $u$ and $v$. If a second pair $\{u, w\}$ with
$w \neq v$ were also infeasible with witness $W_{uw}$, then $W_{uv} \cap W_{uw}$ would
contain only the single terminal~$u$ (so $\delta(W_{uv} \cap W_{uw}) \le 0$), while
$W_{uv} \cup W_{uw}$ would contain the three terminals $u, v, w$ (so
$\delta(W_{uv} \cup W_{uw}) \le 1$). However, $|N_J(\cdot)|$ is submodular, so $\delta$
is supermodular:
\[
  \delta(W_{uv} \cup W_{uw}) + \delta(W_{uv} \cap W_{uw})
  \ge \delta(W_{uv}) + \delta(W_{uw}) \ge 2,
\]
forcing $\delta(W_{uv} \cup W_{uw}) \ge 2$, a contradiction. Hence the infeasible
omitted pairs form a matching on the five terminals, of size at most~$2$. Therefore
$t \in \{8, 9, 10\}$, giving $|\Pi(P)| \in \{36, 45, 55\}$. Finally, at most two of the
five adjacent pairs of a pentagon can be infeasible, leaving at least three feasible
adjacent pairs and hence at least $\binom{4}{2} = 6$ common types with any pentagon.
\end{proof}

\begin{corollary}
Neither a bipartite five-pole nor a one-edge-defect five-pole with all five terminals in
the defect-free part can appear on either side of a cyclic $5$-edge cut of a minimum
counterexample to the Berge--Fulkerson conjecture.
\end{corollary}

\begin{proof}
By Theorem~\ref{thm:replay}, every non-pentagon profile in cases I--XII has size at
least~$25$ and is disjoint from a pentagon, while the two profiles of case~XIII have
sizes $30$ and $26$. By Theorem~\ref{thm:bipartite-21}, a bipartite five-pole has $21$
types (and is not a pentagon), ruling it out. By Theorem~\ref{thm:one-defect}, a
one-edge-defect five-pole has at least $36$ types (strictly more than either side of
case~XIII) and always shares at least $6$ types with every pentagon orientation (ruling
out cases I--XII).
\end{proof}

\begin{thebibliography}{9}
\bibitem{Fulkerson1971}
D.~R.~Fulkerson,
\emph{Blocking and anti-blocking pairs of polyhedra},
Math.\ Program.\ \textbf{1} (1971), 168--194.

\bibitem{MMT}
E.~M\'{a}\v{c}ajov\'{a}, G.~Mazzuoccolo, and G.~Tabarelli,
Ars Math.\ Contemp., \href{https://doi.org/10.26493/1855-3974.3409.c13}{doi:10.26493/1855-3974.3409.c13}.
\end{thebibliography}

\end{document}
